\section{Safeguard Ablations with Live Recovery}
\label{app:live-recovery}

These follow-ups ask whether compiled prefixes save model work on unused
records, whether individual safeguards change deterministic outcomes, and
whether those differences survive a live agent's recovery. Cohorts and endpoints
remain separate from the six studies' 810 compiled episodes in
\cref{sec:discussion}; none changes that historical denominator or supplies a
new end-to-end risk certificate.

\paragraph{Fresh-record comparison.}
On 120 previously unused issues from the same historical snapshot, four live
arms complete all 480 episodes. Baseline, GAC, and support-only each pass the
registered four-field factuality endpoint on 120/120 records; recurrence-only
passes 118/120. On the supplementary seven-field endpoint the respective counts
are 119, 120, 119, and 118. GAC halves model requests (4 to 2 per episode) and
reduces estimated model cost by 45.8\% against baseline. Its dispatch decisions
are identical to support-only: 119 dispatches and one fallback, so separate
stochastic continuations cannot establish a gate effect. The two recurrence-only
misses are strict formatting/case differences, not demonstrated semantic false
assertions. Prefix length and context also differ, preventing an isolated
safeguard attribution. The supplementary conservative 95\% interval for GAC
minus baseline is $[-3.59,3.59]$ points on the registered endpoint; equivalence
and noninferiority are not established. The previously admitted artifact is
fixed, without new discovery or calibration; these are not time-forward
records (\path{experiments/fresh_safeguards/analysis/}).

\paragraph{Deterministic mechanism diagnostic.}
A disjoint cohort supplies 24 discovery, 96 development, 192 calibration, and
96 held-out issue groups. The study specifies a three-comment prefetch, a
requested-limit entry hull, and a response-completeness check; production
helpers fit argument bindings and the gate, but end-to-end mining of these
contracts is not tested. Each held-out issue receives requested limits 1, 3,
and 5 under ordinary responses and responses capped at two comments. Five arms
(full, baseline, and removal of scope, completeness, or gate) yield 2{,}880
provider-free evaluations. Under ordinary responses, full is correct on 288/288
cases versus 260/288 without scope admission; under the cap, full is 288/288
versus 252/288 without completeness checking. The baseline reference policy is
correct throughout. The verifier compensates for scope removal under the
cap, where that ablation remains 288/288. The fitted gate score is constant
(0.00247); all eleven thresholds admit the same 192 calibration groups, and
removing it changes no held-out decision. These are deterministic outcomes,
not live accuracy or evidence that recovery is impossible
(\path{experiments/safeguard_mechanisms/}).

\paragraph{Stopped pilot and separately sealed completion.}
The initial 24-issue live pilot attempted 153 of 576 scheduled episodes:
152 completed, 151 were correct, one raised an execution error, and 423 remained
unattempted. A separate diagnostic returned the correct answer but reproduced
an accounting failure: the old validator rejected reported cache-write tokens.
The original failure retained no usage, so this does not identify its exact
cause. The completion study prices cache writes at the existing frozen rate and
preserves partial usage and unknown-cost reserves. It retains the pilot's
prompt, model, tools, answer contract, prefetch, guards, source, timeout, and
retry policy; its new protocol and schedule were sealed before execution.
Pilot and diagnostic outcomes are excluded from the completion comparison.

\paragraph{Live design and statistical unit.}
The completion uses all 96 offline held-out issues: 72 \emph{confirmation}
groups untouched by the live pilot or diagnostic, and 24 \emph{replication}
groups selected for the pilot, including issues it never reached. Every issue
has prior offline exposure; replication is not independent of the pilot.
Three limits $\times$ two response conditions $\times$ four arms produce 24
episodes per issue, or 2{,}304 overall. Full, no-scope-guard, and
no-completeness-check differ by one intervention; baseline receives no prefetch.
Every arm has the same live continuation and recovery tools, including original
call arguments and completeness information. Tool reads use the pinned local
snapshot, not a remote GitHub service. The constant gate is retained without a
redundant fifth live arm.

The model is \code{gpt-5.6-luna}, with low reasoning and verbosity,
\code{store=false}, no parallel tool calls, at most 2{,}048 output tokens per
request and eight turns, a 120\,s request timeout, and no SDK or transport
retries. Four workers execute issue blocks; all 24 arm permutations occur three
times in confirmation and once in replication within each condition/limit.
All scheduled episodes were attempted, with no replacements: 2{,}300 completed,
2{,}186 were correct, 114 completed with incorrect answers, and four timed out.
Failures count as unsuccessful scheduled attempts (\cref{tab:recovery-success}).

The prespecified paired analysis averages the three limit variants within each
issue, then bootstraps issues 10{,}000 times (seed 20261018), separately for each
cohort and response condition. The 95\% percentile intervals in
\cref{tab:recovery-contrasts} are descriptive and unadjusted for multiple
contrasts. They assume suitable independence between issue groups; author or
repository dependence remains possible. The 2{,}304 episodes are not independent
samples. No noninferiority margin or new live risk certificate is asserted.

\begin{table}[t]
\centering
\caption{Live completion: correct answers / scheduled episodes, including
failures. No scope removes the requested-limit hull; no completeness removes
only that check. All four timeouts occur in confirmation: one each in ordinary
baseline, full, and no scope, plus one in capped full.}
\label{tab:recovery-success}
\small
\setlength{\tabcolsep}{5pt}
\begin{tabular}{@{}llrrrr@{}}
\toprule
Cohort & Responses & Baseline & Full & No scope & No completeness \\
\midrule
Confirmation & Ordinary & 204/216 & 205/216 & 186/216 & 210/216 \\
Confirmation & Capped & 211/216 & 205/216 & 207/216 & 204/216 \\
Replication & Ordinary & 71/72 & 70/72 & 67/72 & 70/72 \\
Replication & Capped & 70/72 & 68/72 & 69/72 & 69/72 \\
\bottomrule
\end{tabular}
\end{table}

\begin{table}[t]
\centering
\caption{Confirmation contrasts: full minus comparator, in percentage points.
Discordances count full-only / comparator-only correct cases. Intervals resample
72 issue groups, not 216 independent requests, and are descriptive/unadjusted.}
\label{tab:recovery-contrasts}
\small
\begin{tabular}{@{}llrrl@{}}
\toprule
Responses & Comparator & Discordances & Difference & 95\% interval \\
\midrule
Ordinary & No scope & 24 / 5 & $+8.80$ & $[3.70,13.89]$ \\
Ordinary & No completeness & 4 / 9 & $-2.31$ & $[-5.56,0.93]$ \\
Ordinary & Baseline & 11 / 10 & $+0.46$ & $[-3.70,4.63]$ \\
Capped & No scope & 6 / 8 & $-0.93$ & $[-4.17,2.31]$ \\
Capped & No completeness & 8 / 7 & $+0.46$ & $[-3.70,4.63]$ \\
Capped & Baseline & 4 / 10 & $-2.78$ & $[-6.48,0.46]$ \\
\bottomrule
\end{tabular}
\end{table}

\paragraph{Scope protection and recoverability.}
Under ordinary responses, scope admission improves confirmation success by
8.80 points. Full-only correct cases occur in 20 issue groups and ablation-only
correct cases in five. Replication gives $+4.17$ points with interval
$[-1.39,9.72]$: directionally consistent but not decisive. A post-hoc sensitivity
on the 214 ordinary pairs where both arms complete gives 23 full-only versus
four no-scope-only correct cases; timeout imbalance does not create the net
19-case advantage, but this does not replace the primary analysis.

Incomplete accepted prefetches are defined by source positions and requested
scope before inspecting live answers. Without scope admission, only 5/20 such
ordinary confirmation cases trigger more reads and only four finish correctly;
in replication, 4/8 read further and finish correctly. Without completeness
checking under the cap, all 27 confirmation cases read further and 22 finish
correctly; all nine replication cases read further and eight finish correctly.
The agent often repairs capped evidence, but frequently treats a page complete
for a three-comment call as sufficient for a larger request. Requests requiring
more than the two-comment cap number 54 in confirmation and 18 in replication.

Full versus no completeness under ordinary responses has identical initial
model inputs: differing accuracy reflects separate stochastic continuations,
not a verifier effect. Under capped responses, completeness can compensate for
the removed scope guard. Neither the unresolved completeness contrast nor the
earlier deterministic gap establishes that every safeguard is necessary.
Correct prefetches also leave semantic errors after full dispatch; accepted
dispatch is not an end-to-end answer certificate.

\paragraph{Resource trade-offs and accounting.}
Every confirmation contrast involving full has an execution failure; therefore
its completed-only resource means are descriptive, and overall saving intervals
are withheld. Replication has complete accounting (\cref{tab:recovery-resources}).
Model requests and estimated model cost fall, while total tools increase and
slightly fewer answers are correct. There is no demonstrated quality equivalence
or overall system-cost win. All discarded prefetch and continuation reads count.
Local snapshot reads, concurrency, caching, frozen rates, and omitted
engineering/human costs limit transfer to deployment.

\begin{table}[t]
\centering
\caption{Complete 24-group replication accounting, baseline $\to$ full.
Model-cost reductions exclude tool-service and lifecycle costs.}
\label{tab:recovery-resources}
\small
\begin{tabular}{@{}lrrrr@{}}
\toprule
Responses & Model requests & Model cost & Total tools & Correct / 72 \\
\midrule
Ordinary & $2.417\to1.931$ & $-18.2\%$ & $1.417\to1.597$ & $71\to70$ \\
Capped & $2.722\to2.458$ & $-9.1\%$ & $1.722\to2.125$ & $70\to68$ \\
\bottomrule
\end{tabular}
\end{table}

The paired model-cost difference intervals are
$[-0.00010541,-0.00005827]$ dollars per episode for ordinary responses and
$[-0.00007654,-0.00001820]$ for capped responses. New-run recorded cost,
including known partial usage, is \$0.952696; four failed-attempt reserves and
the carried pilot/diagnostic accounting give \$7.001762 in estimates plus
reserves. These are neither invoices nor bounds on actual billing. All 448{,}212
cache-write tokens were priced; the prior cache-write rejection did not recur.

\paragraph{Interpretation and provenance.}
The live evidence supports request-scope admission on this constructed workload;
it does not establish a completeness-check accuracy benefit, learned-gate value,
or superiority to baseline. It does not test end-to-end contract mining or
transfer of this guard benefit across models, natural requests, or time. The
broader paper's cross-model and post-cutoff studies concern other questions.
Source regrading is consistency checking, not independent human gold attestation.
All incorrect answers, failure stages, source expectations, metric totals, and
paired intervals are retained in
\path{experiments/live_recovery_completion/analysis.json}; the checkpoint,
prospective seal, diagnostic, and protocol remain unchanged. The post-run audit
regrades answers, checks tool-output digests, and recomputes costs and outcomes.
\Cref{tab:claims-evidence} indexes the drivers and preceding studies.
